\subsection{Entry Admission and Event Depth}
The event-entry state is usable only while (:ref:base.control_registers.ECR.V:) is set. A maskable interrupt is admitted only when (:ref:base.control_registers.ECR.V:),
(:ref:base.registers.STATE.STATUS.IE:), and the depth limit all permit it: the current event depth must be less than (:ref:base.control_registers.ECR.MAX_EDEPTH:). The
interrupt file additionally requires a pending and enabled identity that passes ITHRESH; ITOP chooses the numerically
lowest such identity. Otherwise the identity remains pending in the interrupt file. A synchronous exception that requires delivery while (:ref:base.control_registers.ECR.V:) is
clear cannot be deferred and enters SHUTDOWN. An NMI observed while (:ref:base.control_registers.ECR.V:) is clear sets the hardware-managed
(:ref:base.control_registers.ECR.NMI_P:) latch and is not admitted. After (:ref:base.control_registers.ECR.V:) becomes one, hardware admits that pending NMI before the next
instruction boundary when (:ref:base.registers.STATE.STATUS.NI:) is clear. While (:ref:base.registers.STATE.STATUS.NI:) remains set, the NMI remains pending. Multiple NMI
observations coalesce in (:ref:base.control_registers.ECR.NMI_P:).

Event depth is zero outside event handling. A successful entry saves STATUS in the selected return context and
increments (:ref:base.registers.STATE.STATUS.EDEPTH:). (:ref:base.control_registers.ECR.MAX_EDEPTH:) equal to zero prevents maskable interrupt admission; values 1 through 15 limit
only maskable interrupts and do not suppress synchronous exceptions or NMI. Depth 15 is representable but terminal:
any further event requirement enters SHUTDOWN. At depth 14 or less, a delivery failure can still create an
\texttt{EVENT\_DELIVERY\_FAILURE} family frame; at depth 15 there is no representable child frame.
